\section{Introduction}

% Starte generelt med å utforske hva den faglige diskusjonen går i for tiden
% Gå over til mer spesifikt, hva har andre gjort med problemstillingen som utforskes
% Presenter det vitenskapelige målet og problemstillingen, hva er det de andre ikke har fått til som vi skal gjøre

The separation of microscopic particles according to their physical properties is a central problem in fields such as biophysics, chemistry, and materials science. In particular, the ability to sort biological macromolecules such as DNA by size is essential for analysis and characterization of complex samples. Traditional techniques, such as centrifugation, rely on external forces and macroscopic gradients, but alternative methods based on stochastic transport have attracted significant interest.
\\ \\
One such approach exploits Brownian ratchets, where thermal fluctuations are rectified into directed motion using an asymmetric, time-dependent potential. Although the motion of particles in solution is inherently random due to thermal agitation, it is possible to bias this motion by periodically switching (“flashing”) an asymmetric potential. This leads to a net drift, even in the absence of a net external force, a phenomenon often referred to as a Brownian motor.
\\ \\
In this work, we investigate a model system consisting of particles undergoing Brownian motion in a one-dimensional flashing ratchet potential. The aim is to study how the resulting drift velocity depends on particle size and on the flashing period of the potential, and to assess the feasibility of using this mechanism for particle separation. The problem is approached numerically using a stochastic differential equation framework based on the Langevin equation.